\section*{第 \S\arabic{exercisegroup} 节习题}
\phantomsection
\addcontentsline{toc}{section}{第 \protect\S\arabic{exercisegroup} 节习题}

\begin{enumerate}
\def\labelenumi{\arabic{enumi})}
\tightlist
\item
  设 \((f_{\alpha})\) 是一族实函数，定义在区间 \(\mathrm{I} \subset \mathbf{R}\) 上，其中每个函数在 \(\mathrm{I}\) 上有一个严格原函数，并且它们关于关系 \(\leqslant\) 构成一个有向集。设 \(f\) 是族 \((f_{\alpha})\) 的上包络；假设 \(f\) 在 \(\mathrm{I}\) 上有一个严格原函数。证明：若 \(g_{\alpha}\)（相应地，\(g\)）是 \(f_{\alpha}\)（相应地，\(f\)）的在点 \(x_0 \in \mathrm{I}\) 处为零的原函数，则 \(g\) 是族 \((g_{\alpha})\) 在交集 \(\mathrm{I} \cap [x_0, +\infty[\) 上的上包络，而是它在交集 \(\mathrm{I} \cap ] - \infty\) , \(x_0]\) 上的下包络。（限于这两个区间中的第一个，证明：若 \(u\) 是 \((g_{\alpha})\) 在 \(\mathrm{I} \cap [x_0, +\infty[\) 上的上包络，则 \(u(x + h) - u(x) \leqslant g(x + h) - g(x)\)（对 \(h > 0\)）成立；然后证明，对一切 \(\alpha\)，有 \(u(x + h) - u(x) \geqslant g_{\alpha}(x + h) - g_{\alpha}(x)\)，并由最后这个不等式断定 \(\liminf_{h \to 0} (u(x + h) - u(x))/h \geqslant f(x)\)；由此推出命题。）
\end{enumerate}

给出一个区间 I 上的递增连续函数序列 \((f_n)\) 的例子，这些函数一致有界，但其上包络在 I 上没有严格原函数。

\begin{enumerate}
\def\labelenumi{\arabic{enumi})}
\setcounter{enumi}{1}
\item
  证明：为使函数 \(\mathbf{f}\) 是区间 I 上的阶梯函数，必须且只需它只有有限个不连续点，并且在它连续的每个区间上取常值。
\item
  设 \(\mathbf{f}\) 是区间 \(I \subset \mathbf{R}\) 上的规则函数，取值于完备赋范空间 \(E\)（\(\mathbf{R}\) 上的）中；证明：对每个紧子集 \(H\)（含于 \(I\)），集合 \(\mathbf{f}(H)\) 在 \(E\) 中是相对紧的；给出一个 \(\mathbf{f}(H)\) 在 \(E\) 中不是闭集的例子。
\item
  给出一个连续实函数 \(f\)（在紧区间 \(\mathrm{I} \subset \mathbf{R}\) 上）的例子，使得复合函数 \(x \mapsto \operatorname{sgn}(f(x))\) 在 \(\mathrm{I}\) 上不是规则的（尽管 \(\operatorname{sgn}\) 在 \(\mathbf{R}\) 上是规则的）。
\item
  设 \(\mathbf{f}\) 是定义在紧区间 \(\mathrm{I} = [a, b] \subset \mathbf{R}\) 上、取值于完备赋范空间 \(\mathrm{E}\) 中的向量函数；我们说 \(\mathbf{f}\) 在 \(\mathrm{I}\) 上是有界变差的，如果存在一个数 \(m > 0\)，使得对每个有限的严格递增序列 \((x_i)_{0 \leqslant i \leqslant n}\)——由
\end{enumerate}

I 中满足 \(x_0 = a\) 与 \(x_n = b\) 的点组成——有 \(\sum_{i=0}^{n-1} \| \mathbf{f}(x_{i+1}) - \mathbf{f}(x_i) \| \leqslant m\)。

\begin{enumerate}
\def\labelenumi{\alph{enumi})}
\item
  证明 \(\mathbf{f}(\mathrm{I})\) 在 E 中相对紧（用反证法论证）。
\item
  证明 \(\mathbf{f}\) 在 I 上是规则的（证明当 \(x\) 趋于一点 \(x_0 \in \mathrm{I}\) 而保持 \(>x_0\) 时，函数 \(\mathbf{f}\) 不可能有两个不同的聚点，并利用 \(a\)））。
\end{enumerate}

\begin{enumerate}
\def\labelenumi{\arabic{enumi})}
\setcounter{enumi}{5}
\tightlist
\item
  为使函数 \(\mathbf{f}\)（取值于完备赋范空间 \(E\) 中、定义在开区间 \(I \subset R\) 上）在 \(I\) 的一个可数子集的余集的所有点处等于 \(I\) 上的一个规则函数，必须且只需它满足下列条件：对每个 \(x \in I\) 与每个 \(\varepsilon > 0\)，存在一个数 \(h > 0\) 和两个元素 \(a\)、\(b\)（\(E\) 的），使得 \(\| \mathbf{f}(y) - a \| \leqslant \varepsilon\) 对每个 \(y \in [x, x + h]\) 都成立（至多除去这个区间的可数无穷多个点），且 \(\| \mathbf{f}(z) - b \| \leqslant \varepsilon\) 对所有 \(z \in [x - h, x]\) 成立（至多除去这个区间的可数无穷多个点）。
\end{enumerate}

* 7) 证明：等于 \(\sin(1/x)\)（当 \(x \neq 0\)）而等于 0（当 \(x = 0\)）的函数在 \(\mathbf{R}\) 上有一个严格原函数（注意 \(x^2 \sin(1/x)\) 在每一点处有导数）。

由此推出：若 \(g(x,u,v)\) 是 u, v 的多项式，其系数是 x 在区间 I（包含 0）上的连续函数，则等于 \(g(x,\sin1/x,\cos1/x)\)（当 \(x\neq0\)）而等于一个适当的值 \(\alpha\)（待定）（当 x=0）的函数在 I 上有一个严格原函数；给出一个 \(\alpha\neq g(0,0,0)\) 的例子。

* 8) 证明：存在 \([-1, +1]\) 上的连续函数，在这个区间的每一点处有有限导数，这个导数等于 \(\sin\left(\frac{1}{\sin 1/x}\right)\)（在每一点 \(x\) 处，只要 x 异于 \(1/n\pi\)（\(n\) 为整数 \(\neq 0\)）且异于 0）。（在 \(x = 1/n\pi\) 的邻域内作变量替换 \(x = \frac{1}{n\pi + \operatorname{Arc} \sin u}\) 并利用习题 7；借助同样的变量替换证明，存在一个常数 \(a > 0\)（与 \(n\) 无关），使得

\[
\left| \int_ {2 / (2 n + 1) \pi} ^ {2 / (2 n - 1) \pi} \sin \left(\frac {1}{\sin \frac {1}{x}}\right) d x \right| \leqslant \frac {a}{n ^ {3}};
\]

并由此断定，若令

\[
g (x) = \lim _ {\varepsilon \rightarrow 0} \int_ {\varepsilon} ^ {x} \sin \left(\frac {1}{\sin \frac {1}{t}}\right) d t,
\]

则 g 在点 x = 0 处有一个等于 0 的导数。）*

\begin{enumerate}
\def\labelenumi{\arabic{enumi})}
\setcounter{enumi}{8}
\item
  设 \(\mathbf{f}\) 是紧区间 \(\mathrm{I} = [a, b]\) 上的规则函数。证明：对每个 \(\varepsilon > 0\)，存在一个连续函数 \(\mathbf{g}\)（在 \(\mathrm{I}\) 上），使得 \(\int_{a}^{b} \| \mathbf{f}(t) - \mathbf{g}(t) \| dt \leqslant \varepsilon\)（化为 \(\mathbf{f}\) 是阶梯函数的情形）。由此推出存在一个多项式 \(\mathbf{h}\)（系数在 E 中），使得 \(\int_{a}^{b} \| \mathbf{f}(t) - \mathbf{h}(t) \| dt \leqslant \varepsilon\)。
\item
  设 \(\mathbf{f}\) 是 \([a, b]\) 上的规则函数，取值于 E 中；设 \(\mathbf{g}\) 是 \([a, c]\)（\(c > b\)）上的规则函数，取值于 F 中；并设 \((\mathbf{x}, \mathbf{y}) \mapsto [\mathbf{x}, \mathbf{y}]\) 是 \(\mathrm{E} \times \mathrm{F}\) 到 G 中的连续双线性映射（E, F, G 是完备赋范空间）。证明
\end{enumerate}

\[
\lim _ {h \rightarrow 0, h > 0} \int_ {a} ^ {b} [ \mathbf {f} (t). \mathbf {g} (t + h) ] d t = \int_ {a} ^ {b} [ \mathbf {f} (t). \mathbf {g} (t) ] d t
\]

（化为 \(\mathbf{f}\) 是阶梯函数的情形）。

\begin{enumerate}
\def\labelenumi{\arabic{enumi})}
\setcounter{enumi}{10}
\tightlist
\item
  在与习题 10 相同的假设下，证明：对每个 \(\varepsilon > 0\)，存在一个数 \(\rho > 0\)，使得把 \([a, b]\) 分成区间 \([x_i, x_{i+1}]\)（长度 \(\leqslant \rho\)，\(0 \leqslant i \leqslant n-1\)）的每个分割，有
\end{enumerate}

\[
\left\| \int_ {a} ^ {b} [ \mathbf {f} (t). \mathbf {g} (t) ] d t - \sum_ {i = 0} ^ {n - 1} [ \mathbf {f} (u _ {i}). \mathbf {g} (v _ {i}) ] (x _ {i + 1} - x _ {i}) \right\| \leqslant \varepsilon
\]

对每个指标 i，不论点 \(u_{i}\)、\(v_{i}\)（取自 \([x_{i}, x_{i+1}]\)）如何选取（化为 f 与 g 都是阶梯函数的情形）。

¶ 12) 我们说实数序列 \((x_{n})\)（含于区间 \([0, 1]\)）在这个区间中是一致分布的，如果

\[
\lim _ {n \to \infty} \frac {\nu_ {n} (\alpha , \beta)}{n} = \beta - \alpha\tag{1}
\]

对每对数 \(\alpha\)、\(\beta\)（满足 \(0 \leqslant \alpha \leqslant \beta \leqslant 1\)）都成立，其中 \(\nu_{n}(\alpha, \beta)\) 表示指标 \(i\)（满足 \(1 \leqslant i \leqslant n\) 且 \(\alpha \leqslant x_{i} \leqslant \beta\)）的个数。

证明：若序列 \((x_{n})\) 一致分布且 \(\mathbf{f}\) 是 \([0, 1]\) 上的规则函数，则

\[
\lim _ {n \rightarrow \infty} \frac {1}{n} \sum_ {i = 1} ^ {n} \mathbf {f} (x _ {i}) = \int_ {0} ^ {1} \mathbf {f} (t) d t\tag{2}
\]

（化为 \(\mathbf{f}\) 是阶梯函数的情形）。逆命题。

证明：为使序列 \((x_{n})\) 一致分布，只需关系 (2) 对每个实函数 f（属于赋以一致收敛拓扑的 {[}0, 1{]} 上实连续函数空间中的一个稠密集）成立。

¶ 13) 设 \(f\) 是紧区间 \([a, b]\) 上的实规则函数。令

\[
r (n) = \frac {b - a}{n} \sum_ {k = 1} ^ {n} f \left(a + k \frac {b - a}{n}\right) - \int_ {a} ^ {b} f (t) d t.
\]

\begin{enumerate}
\def\labelenumi{\alph{enumi})}
\tightlist
\item
  证明：若 \(f\) 在 \([a, b]\) 上递增，则
\end{enumerate}

\[
0 \leqslant r (n) \leqslant \frac {b - a}{n} \Bigl (f (b) - f (a) \Bigr).
\]

\begin{enumerate}
\def\labelenumi{\alph{enumi})}
\setcounter{enumi}{1}
\tightlist
\item
  若 \(f\) 连续且在 \([a, b[\) 上有一个有界的规则右导数，证明有 \(\lim_{n\to \infty}nr(n) = \frac{b - a}{2} (f(b) - f(a))\)（令 \(x_{k} = a + k\frac{b - a}{n}\)），并注意有
\end{enumerate}

\[
r (n) = \sum_ {k = 0} ^ {n - 1} \int_ {x _ {k}} ^ {x _ {k + 1}} \left(f (x _ {k + 1}) - f (t)\right) d t,
\]

并应用 II, p.~57 的命题 5）。

\begin{enumerate}
\def\labelenumi{\alph{enumi})}
\setcounter{enumi}{2}
\tightlist
\item
  给出一个函数 \(f\) 的例子，它在 \([a, b]\) 上递增且连续，使得 \(nr(n)\) 不趋于 \(\frac{b - a}{2} (f(b) - f(a))\)（当 \(n\) 无限增大时） {[}取 \(f\) 为递增函数的递减序列 \((f_n)\) 的极限，这些递增函数的图象是折线，使得
\end{enumerate}

\[
f _ {n} \left(a + k \frac {b - a}{2 ^ {n}}\right) = f \left(a + k \frac {b - a}{2 ^ {n}}\right) \quad \text { for } \quad 0 \leqslant k \leqslant 2 ^ {n}
\]

且

\[
(b - a) \sum_ {k = 1} ^ {2 ^ {n}} f _ {n} \left(a + k \frac {b - a}{2 ^ {n}}\right) - 2 ^ {n} \int_ {a} ^ {b} f _ {n} (t) d t \geqslant \frac {3}{4} (b - a) \left(f _ {n} (b) - f _ {n} (a)\right) \Biggr ].
\]

\begin{enumerate}
\def\labelenumi{\arabic{enumi})}
\setcounter{enumi}{13}
\tightlist
\item
  设 \(\mathbf{f}\) 是一个向量函数，它是某个规则函数 \(\mathbf{f}'\)（在 \([a, b]\) 上）的原函数，且满足 \(\mathbf{f}(a) = \mathbf{f}(b) = 0\)。证明：若 \(\mathbf{M}\) 是 \(\| \mathbf{f}'(x)\|\) 在 \([a, b]\) 中使 \(\mathbf{f}'\) 连续的点所成之集上的上确界，则
\end{enumerate}

\[
\left\| \int_ {a} ^ {b} \mathbf {f} (t) d t \right\| \leqslant \mathrm{M} \frac {(b - a) ^ {2}}{4}.
\]

\begin{enumerate}
\def\labelenumi{\arabic{enumi})}
\setcounter{enumi}{14}
\tightlist
\item
  设 \(\mathbf{f}\) 是连续实函数，在区间 \([0, a]\) 上严格递增，且满足 \(f(0) = 0\)；设 \(g\) 是它的反函数，定义在 \([0, f(a)]\) 上且严格递增；证明
\end{enumerate}

\[
x y \leqslant \int_ {0} ^ {x} f (t) d t + \int_ {0} ^ {y} g (u) d u
\]

对 \(0 \leqslant x \leqslant a\) 与 \(0 \leqslant y \leqslant f(a)\) 成立，且仅当 \(y = f(x)\) 时取等号（把 \(xy - \int_{0}^{x} f(t) dt\) 作为 x 的函数研究其变差（y 固定））。由此推出：对 \(x \geqslant 0\)，\(y \geqslant 0\)，p \textgreater{} 1，\(p' = p/(p - 1)\)，有 \(xy \leqslant ax^{p} + by^{p'}\)（对 a \textgreater{} 0，b \textgreater{} 0 且 \((pa)^{p'}(p'b)^{p} \geqslant 1\)）。

\begin{enumerate}
\def\labelenumi{\arabic{enumi})}
\setcounter{enumi}{15}
\tightlist
\item
  设 \(\mathbf{f}\) 是 \(\mathrm{I} = [a, b] \subset \mathbf{R}\) 上的规则向量函数，\(\mathbf{u}\) 是 \(\mathbf{f}\) 在 \(\mathrm{I}\) 上的一个原函数，\(\mathrm{D}\) 是包含 \(\mathbf{u}(\mathrm{I})\) 的闭凸集。证明：若 \(g\) 是 \(\mathrm{I}\) 上的实单调函数，则
\end{enumerate}

\[
\int_ {a} ^ {b} \mathbf {f} (t) g (t) d t = (\mathbf {u} (b) - \mathbf {c}) g (b) + (\mathbf {c} - \mathbf {u} (a)) g (a)
\]

其中 \(\mathbf{c}\) 属于 D（化为 \(g\) 是单调阶梯函数的情形）。由此推出：若 \(f\) 是 I 上的实规则函数，则存在一个 \(c \in \mathrm{I}\)，使得

\[
\int_ {a} ^ {b} f (t) g (t) d t = g (a) \int_ {a} ^ {c} f (t) d t + g (b) \int_ {c} ^ {b} f (t) d t
\]

（“第二中值定理”）。

\begin{enumerate}
\def\labelenumi{\arabic{enumi})}
\setcounter{enumi}{16}
\tightlist
\item
  设 \(g\) 是实函数，有连续导数，且 \(\neq 0\)（在 \([a, x]\) 上）；若 \(f\) 是实函数，具有规则的 \((n + 1)^{th}\) 导数（在 \([a, x]\) 上），证明余项 \(r_n(x)\)——即 \(n\) 阶泰勒展开（\(f\) 在点 \(a\) 处的）中的余项——可以写成
\end{enumerate}

\[
r _ {n} (x) = (g (x) - g (a)) \frac {(x - \xi) ^ {n}}{n !} \frac {f ^ {(n + 1)} (\xi)}{g ^ {\prime} (\xi)}
\]

其中 \(a < \xi < x\)（利用 \(r_{n}(x)\) 的积分形式）。

\begin{enumerate}
\def\labelenumi{\arabic{enumi})}
\setcounter{enumi}{17}
\tightlist
\item
  设 \(f\) 是有限实函数，在开区间 \(I\) 上连续。为使 \(f\) 在 \(I\) 上是凸的，必须且只需对每个 \(x \in I\) 有
\end{enumerate}

\[
\operatorname * {l i m s u p} _ {h \to \infty} \left(\frac {1}{2 h} \int_ {x - h} ^ {x + h} f (t) d t - f (x)\right) \geqslant 0
\]

（按 I, p.~46 习题 9 的方式论证）。

\begin{enumerate}
\def\labelenumi{\arabic{enumi})}
\setcounter{enumi}{18}
\tightlist
\item
  设 \(f\) 是区间 \(I\) 上的凸函数，\(h\) 是一个数 \(> 0\)，\(I_h\) 是 \(I\) 与区间 \(I + h\)、\(I - h\) 的交；证明：若 \(I_h\) 非空，则函数
\end{enumerate}

\[
g _ {h} (x) = \frac {1}{2 h} \int_ {x - h} ^ {x + h} f (t) d t
\]

在 \(\mathrm{I}_h\) 上是凸的；若 \(h < k\)，则 \(g_h \leqslant g_k\)。证明：当 \(h\) 趋于 0 时，\(g_h\) 一致收敛于 \(f\)（在含于 \(\mathrm{I}\) 的内部的每个紧区间上）。

\begin{enumerate}
\def\labelenumi{\arabic{enumi})}
\setcounter{enumi}{19}
\tightlist
\item
  证明：当 \(n\) 无限增大时，多项式
\end{enumerate}

\[
f _ {n} (x) = \frac {\int_ {0} ^ {x} (1 - t ^ {2}) ^ {n} d t}{\int_ {0} ^ {1} (1 - t ^ {2}) ^ {n} d t}
\]

在每个区间 \([-1, -\varepsilon]\) 上一致收敛于 -1，而在每个区间 \([\varepsilon, +1]\) 上一致收敛于 +1，其中 \(\varepsilon > 0\)（注意 \(\int_{0}^{1}(1-t^{2})^{n}dt \geqslant \int_{0}^{1}(1-t)^{n}dt\)）。由此推出：多项式 \(g_{n}(x) = \int_{0}^{x} f_{n}(t)dt\) 一致收敛于 \(|x|\)（在 \([-1, +1]\) 上），这给出了魏尔斯特拉斯定理的一个新证明（Gen.~Top., X, p.~313, 命题 3）。

\begin{enumerate}
\def\labelenumi{\arabic{enumi})}
\setcounter{enumi}{20}
\tightlist
\item
  设 f 是区间 \([0, a]\) 上的实递增凸连续函数，且满足 \(f(0) = 0\)。若 \(a_{1} \geqslant a_{2} \geqslant \cdots \geqslant a_{n} \geqslant 0\) 是 \([0, a]\) 中点的一个有限递减序列，证明有
\end{enumerate}

\[
f \left(a _ {1}\right) - f \left(a _ {2}\right) + \dots + (- 1) ^ {n - 1} f \left(a _ {n}\right) \geqslant f \left(a _ {1} - a _ {2} + \dots + (- 1) ^ {n - 1} a _ {n}\right).
\]

（可以只限于 \(n = 2m\) 为偶数的情形；注意对 \(1 \leqslant j \leqslant m\) 有

\[
\int_ {\alpha} ^ {\beta} f ^ {\prime} (t) d t \leqslant \int_ {a _ {2 j}} ^ {a _ {2 j - 1}} f ^ {\prime} (t) d t
\]

其中 \(\alpha = a_{2j+1} - a_{2j+2} + \cdots + a_{2m-1} - a_{2m}\)，\(\beta = \alpha + (a_{2j-1} - a_{2j})\)。

\begin{enumerate}
\def\labelenumi{\arabic{enumi})}
\setcounter{enumi}{21}
\tightlist
\item
  设 f 是连续递增的实函数，且 \(\geqslant 0\)（在区间 {[}0, 1{]} 上）。
\end{enumerate}

\begin{enumerate}
\def\labelenumi{\alph{enumi})}
\item
  证明：存在一个凸函数 \(g \geqslant 0\)（在 {[}0,1{]} 上），使得 \(g \leqslant f\) 且 \(\int_0^1 g(t)dt \geqslant \frac{1}{2}\int_0^1 f(t)dt\)（参见 I, p.~48, 习题 23）。
\item
  证明：存在一个凸函数 \(h\)（在 {[}0, 1{]} 上），使得 \(h \geqslant f\)，并且
\end{enumerate}

\[
\int_ {0} ^ {1} h (t) d t \leqslant 2 \int_ {0} ^ {1} f (t) d t.
\]

（对每个 \(a\)（满足 \(0 < a \leqslant 1\)），令 \(f_a\) 是等于 \(f\)（当 \(0 \leqslant t \leqslant a\)）而等于 \(f(a)\)（当 \(a \leqslant t \leqslant 1\)）的函数；设 A 是使下列事实成立的 \(a \in [0,1]\) 所成的集：存在凸函数 \(h_a\)（在 \([0,1]\) 上），满足 \(h_a \geqslant f_a\) 且 \(\int_0^1 h_a(t)dt \leqslant 2\int_0^1 f_a(t)dt\)。证明上确界 \(b\)（A 的）仍属于 A，利用 I, p.~45 的习题 1。然后用反证法证明 \(f_{b} = f\)。为此，归结为证明下列结果：若 \(\varphi\) 在 \([0, 1]\) 上连续且递增、非常数、且满足 \(\varphi(0) = 0\)，则存在一点 \(c \in ]0, 1[\)，使得 \(\varphi(c) > 0\)，并且使得线性函数 \(\psi\)（满足 \(\psi(c) = \varphi(c)\) 及

\[
\psi (2 c - 1) = 0
\]

满足关系式 \(\psi(t) \geqslant \varphi(t)\)（对 \(\max(0, 2c - 1) \leqslant t \leqslant c\)）。

\begin{enumerate}
\def\labelenumi{\arabic{enumi})}
\setcounter{enumi}{22}
\tightlist
\item
  设 \(\mathbf{f}\) 是向量函数，在区间 \([a, b] \subset \mathbf{R}\) 上连续可微，取值于完备赋范空间中。
\end{enumerate}

\begin{enumerate}
\def\labelenumi{\alph{enumi})}
\tightlist
\item
  证明：对 \(a \leqslant t \leqslant b\)，有
\end{enumerate}

\[
\mathbf {f} (t) = \frac {1}{b - a} \int_ {a} ^ {b} \mathbf {f} (x) d x + \int_ {a} ^ {b} \frac {x - a}{b - a} \mathbf {f} ^ {\prime} (x) d x + \int_ {t} ^ {b} \frac {x - b}{b - a} \mathbf {f} ^ {\prime} (x) d x.
\]

\begin{enumerate}
\def\labelenumi{\alph{enumi})}
\setcounter{enumi}{1}
\tightlist
\item
  推出不等式
\end{enumerate}

\[
\| \mathbf {f} (t) \| \leqslant \frac {1}{b - a} \int_ {a} ^ {b} \| \mathbf {f} (x) \| d x + \int_ {a} ^ {b} \left\| \mathbf {f} ^ {\prime} (x) \right\| d x
\]

及

\[
\left\| \mathbf {f} \left(\frac {1}{2} (a + b)\right) \right\| \leqslant \frac {1}{b - a} \int_ {a} ^ {b} \| \mathbf {f} (x) \| d x + \frac {1}{2} \int_ {a} ^ {b} \left\| \mathbf {f} ^ {\prime} (x) \right\| d x.
\]

\setcounter{exercisegroup}{1}
\refstepcounter{exercisegroup}
